A migração e a modernização de sistemas robóticos legados envolvem a adaptação do sistema aos mecanismos de novos \textit{frameworks} e a refatoração da implementação, ao mesmo tempo em que se busca preservar os comportamentos essenciais do sistema original. Nesse contexto, este trabalho propõe uma abordagem sistemática de migração orientada pelo estabelecimento de uma \textit{Ground Truth} comportamental do sistema legado. A abordagem integra a especificação e a verificação de comportamentos, a identificação de suas realizações na implementação legada, a análise das transformações de migração e a validação comportamental do sistema migrado. A proposta foi examinada por meio de um estudo de caso parcialmente retrospectivo da migração da \textit{Body Sensor Network} (BSN) do ROS~1 para o ROS~2. Os resultados indicam a possibilidade de sistematizar retrospectivamente o processo de migração da BSN, estabelecendo relações entre contratos comportamentais, responsabilidades de implementação e transformações de migração. A análise permitiu caracterizar três transformações relacionadas ao gerenciamento do ciclo de vida dos nós, à ampliação do monitoramento e à reorganização das responsabilidades de processamento, comunicação e avaliação de risco. A comparação dos cenários especificados segundo o \textit{Behavior-Driven Development} (BDD), de seus oráculos e dos resultados de execução permitiu identificar evidências de preservação comportamental para os comportamentos selecionados, bem como limitações nos critérios de verificação. O caráter parcialmente retrospectivo do estudo, entretanto, limita as conclusões sobre a eficácia da abordagem na condução do processo original de migração.